% vim: ai sw=2 spell spelllang=cs

\begin{frame}
  \frametitle{Paměť jinak I.}

  \heading{LDD3 kap. 15 (zastaralá)}

  \iffimuni {
    \bigskip

    \tableofcontents

    \heading{Příště}
    \begin{itemize}
      \item Přímý přístup do paměti (DMA)
    \end{itemize}
  } \fi

\end{frame}

\mysection{Mapování paměti jádra (\code{mmap})}

\iffimuni
\subsection{\code{mmap} v uživatelském prostoru}
\fi

\begin{frame}[fragile]
  \frametitle{Mapování paměti jádra}

  \heading{Předání dat do/z procesu}
    \begin{itemize}
      \item Známe: \code{read}, \code{write}, \code{ioctl}, \dots

      \pause

      \item U všeho nutné kopírování dat
      \begin{itemize}
	\item \code{copy_\{from,to\}_user} apod.
      \end{itemize}
  \end{itemize}

  \pause
  \bigskip

  \heading{Mapování paměti}
  \begin{itemize}
    \item Namapování stránek do procesu
    \item Proces používá kus stejné paměti jako jádro
  \end{itemize}

  \pause

  \begin{block}{Systémové volání \code{mmap} (uživatelský prostor)}
  \begin{lstlisting}
void *mmap(void *addr, size_t len, int prot, int flags, int fd, off_t off);
  \end{lstlisting}
  \end{block}

\end{frame}

\begin{frame}
  \frametitle{Úkol}

  \heading{Alokace pomocí \code{mmap} (uživatelský prostor)}
  \begin{enumerate}
    \item Upravujte nějaký \code{main.c} v uživatelském prostoru
    \item Anonymní paměť
    \begin{itemize}
      \item \footnotesize{\code{void *mmap(void *addr, size_t len,
	int prot, int flags, int fd, off_t off)}}
      \item \code{len}: 20M
      \item \code{prot}: \code{PROT_READ | PROT_WRITE}
      \item \code{flags}: \code{MAP_PRIVATE | MAP_ANONYMOUS}
      \item \code{fd}: -1
    \end{itemize}

    \pause

    \item Mapování \file{/dev/zero}
    \begin{itemize}
      \item \code{len} a \code{prot} stejné
      \item \code{flags}: \code{MAP_PRIVATE}
      \item \code{fd}: deskriptor otevřeného \code{/dev/zero}
    \end{itemize}
  \end{enumerate}
\end{frame}

\iffimuni
\subsection{\code{mmap} v jádře}
\fi

\begin{frame}[fragile]{Parametry \code{mmap} v jádře}

  \begin{itemize}
    \item V uživatelském prostoru
    \begin{itemize}
      \item \footnotesize{\code{void *mmap(void *addr, size_t len,
	int prot, int flags, int fd, off_t off)}}
    \end{itemize}

    \pause

    \item V jádře (opět) jedna položka ve \code{struct file_operations}
    \begin{itemize}
      \item \footnotesize{\code{int mmap(struct file *filp,
	struct vm_area_struct *vma)}}

      \pause

      \item Parametry jsou předány přes \code{struct vm_area_struct}

      \onslide<10->

      \item Ovladač musí mapovat stránky mezi \code{vm_start} a \code{vm_end}
      \item Ale \emph{POZOR}, také ověřit privilegia (čtení, zápis, spuštění).

    \end{itemize}
  \end{itemize}

  \onslide<4->
  \pause

  \begin{block}{\code{struct vm_area_struct}}
    \begin{lstlisting}[belowskip=0pt]
  unsigned long vm_start; /* addr or random when addr is NULL */
  unsigned long vm_end;   /* vm_end = vm_start+len */
    \end{lstlisting}
    \pause
    \begin{lstlisting}[aboveskip=0pt,belowskip=0pt]
  unsigned long vm_pgoff; /* vm_pgoff = off/PAGE_SIZE */
    \end{lstlisting}
    \pause
    \begin{lstlisting}[aboveskip=0pt,belowskip=0pt]
  unsigned long vm_flags; /* vm_flags = encoded(flags|prot), see VM_READ etc. */
    \end{lstlisting}
    \pause
    \begin{lstlisting}[aboveskip=0pt,belowskip=0pt]
  pgprot_t vm_page_prot;  /* only for remap_* functions */
  ...
    \end{lstlisting}
    \pause
    \begin{lstlisting}[aboveskip=0pt,belowskip=0pt]
  const struct vm_operations_struct *vm_ops; /* later ... */
    \end{lstlisting}
    \pause
    \begin{lstlisting}[aboveskip=0pt]
  void *vm_private_data;
    \end{lstlisting}
  \end{block}

\end{frame}

\begin{frame}[fragile]{Základní \code{mmap} funkce}

  \begin{itemize}
    \item \header{linux/mm.h}, \code{struct vm_area_struct}
    \item \code{__get_free_page/pages} $\Rightarrow$
      \code{remap_pfn_range}
    \item \code{vmalloc_user} $\Rightarrow$ \code{remap_vmalloc_range}
  \end{itemize}

  \begin{block}{Příklad}
    \begin{lstlisting}[aboveskip=2pt,belowskip=2pt]
int my_init(void)
{
  mem = __get_free_pages(GFP_KERNEL, 2);
  /* mem = vmalloc_user(PAGE_SIZE); */
}
    \end{lstlisting}

    \pause

    \begin{lstlisting}[aboveskip=2pt,belowskip=2pt]
int my_mmap(struct file *filp, struct vm_area_struct *vma)
{
  if ((vma->vm_flags & (VM_WRITE | VM_READ)) != VM_READ)
    return -EINVAL;
  return remap_pfn_range(vma, vma->vm_start, page_to_pfn(virt_to_page(mem)),
      4 * PAGE_SIZE, vma->vm_page_prot);
  /* return remap_vmalloc_range(vma, mem, 0); */
}
    \end{lstlisting}
  \end{block}

\end{frame}

\begin{frame}{Úkol}

  \heading{Přemapování 2 prostorů}
  \begin{enumerate}
    \item Alokujte 2+2 stránky
    \begin{itemize}
      \item Dvojstránku pomocí \code{vmalloc_user}
      \item Dvojstránku pomocí stránkového alokátoru
    \end{itemize}

    \item Zapište na všechny 4 stránky libovolný, ale různý řetězec

    \item Vystavte (mapujte) stránky v \code{mmap}
    \begin{itemize}
      \item První dvojice RO, druhá R/W (ověřte
        \code{prot}, tj.\ \code{vma->vm_flags})
      \item $0 \leq \text{\code{vma->vm_pgoff}} < 2$ $\Rightarrow$ jedno
        mapování
      \item $3 \leq \text{\code{vma->vm_pgoff}} < 4$ $\Rightarrow$ druhé
        mapování
    \end{itemize}

    \item Z userspace vyzkoušejte
    \begin{itemize}
      \item Prostudujte a spusťte \file{pb173/10/pb173.c}
    \end{itemize}
  \end{enumerate}

  \medskip

  \heading{API (opakování)}
  \begin{itemize}
    \item \header{linux/mm.h}, \code{struct vm_area_struct}
    \item \footnotesize{\code{int mmap(struct file *filp,
      struct vm_area_struct *vma)}}
    \item \code{__get_free_page*} $\Rightarrow$ \code{remap_pfn_range}
    \item \code{vmalloc_user} $\Rightarrow$ \code{remap_vmalloc_range}
  \end{itemize}

\end{frame}

\iffimuni
\subsection{\code{mmap} po stránkách}
\fi

\begin{frame}[fragile]{\code{mmap} po stránkách}

  \begin{block}{\code{struct vm_operations_struct}}
    \begin{lstlisting}[aboveskip=2pt,belowskip=2pt]
void (*open)(struct vm_area_struct *vma);
void (*close)(struct vm_area_struct *vma);
int (*fault)(struct vm_area_struct *vma, struct vm_fault *vmf);
    \end{lstlisting}
  \end{block}

  \heading{Přemapování roztroušených stránek}
  \begin{itemize}
    \item Přes \code{remap_pfn_range} obtížně
    \item Při výpadcích stránek se mapují takové stránky jednotlivě
    \begin{itemize}
      \item Každý ovladač má ,,page fault handler''
      \item Háčky \code{struct vm_operations_struct} a v ní \code{fault}
    \end{itemize}

    \pause

    \item V \code{mmap}/\code{fault} je třeba zkontrolovat rozsahy a velikosti
    \begin{itemize}
      \item Předtím to dělaly \code{remap_*_range} funkce
    \end{itemize}

    \pause

    \item \code{vma->vm_ops} se nastaví v \code{mmap}
    \begin{itemize}
      \item Na strukturu s háčky (zejm.\ \code{fault})
    \end{itemize}

    \pause

    \item \code{vma->vm_private_data}
    \begin{itemize}
      \item Pro naše potřeby
      \item K předání informací z \code{file_operations->mmap} do
        \code{vm_ops->*}
    \end{itemize}
  \end{itemize}

\end{frame}

\begin{frame}[fragile]{\code{mmap} po stránkách -- příklad}

  \begin{lstlisting}[belowskip=0pt]
int my_fault(struct vm_area_struct *vma, struct vm_fault *vmf)
{ /* my_data == vma->vm_private_data; */
  \end{lstlisting}
  \begin{uncoverenv}<2->
    \begin{lstlisting}[aboveskip=0pt,belowskip=0pt]
  unsigned long offset = vmf->pgoff << PAGE_SHIFT;
  struct page *page;

  page = my_find_page(offset);
  if (!page)
    return VM_FAULT_SIGBUS;
  get_page(page);
  vmf->page = page;
  return 0;
    \end{lstlisting}
  \end{uncoverenv}
  \begin{lstlisting}[aboveskip=0pt]
}

struct vm_operations_struct my_vm_ops = { .fault = my_fault, };

int my_mmap(struct file *filp, struct vm_area_struct *vma)
{ /* don't forget to check ranges */
  vma->vm_ops = &my_vm_ops;
  vma->vm_private_data = my_data;
  return 0;
}
  \end{lstlisting}

\end{frame}

\begin{frame}{Úkol}

  \heading{Mapování roztroušených stránek} \iffimuni (součást domácího) \fi
  \begin{enumerate}
    \item Předchozí příklad rozšiřte
    \item Místo alokací dvoustránek alokujte 4 samostatné stránky
    \begin{itemize}
      \item 2$\times$ \code{vmalloc_user(PAGE_SIZE)} a 2$\times$ \code{__get_free_page}
    \end{itemize}

    \item Přemapujte stránky v \code{mmap}
    \begin{itemize}
      \item Změna \code{remap_pfn_range} na \code{vma->vm_ops->fault}
    \end{itemize}

    \item Userspace program stejný
    \begin{itemize}
      \item Funkčnost navenek musí být zachovaná
    \end{itemize}
  \end{enumerate}

  \medskip

  \heading{Potřebné podkroky}
  \begin{itemize}
    \item Definice {\footnotesize\code{int fault(struct vm_area_struct *vma,
      struct vm_fault *vmf);}}
    \item Kontrola rozsahů v \code{mmap} (\code{end-start < 2*PAGE_SIZE}
      apod.)
    \item Definice \code{vm_ops} a přiřazení do \code{vma->vm_ops}
  \end{itemize}

\end{frame}
